% ─────────────────────────────────────────────────────────────
% Informe — Taller "SecureChat" · CE-1105 Seguridad de la Información
% Compilar:  pdflatex Informe_SecureChat.tex   (dos veces, por el índice)
% ─────────────────────────────────────────────────────────────
\documentclass[11pt,a4paper]{article}

\usepackage[utf8]{inputenc}
\usepackage[T1]{fontenc}
\usepackage[spanish,es-noshorthands]{babel}
\usepackage[margin=2.5cm]{geometry}
\usepackage{graphicx}
\usepackage{xcolor}
\usepackage{listings}
\usepackage{booktabs}
\usepackage{tabularx}
\usepackage{amssymb}
\usepackage{hyperref}
\hypersetup{colorlinks=true, linkcolor=blue, urlcolor=blue}

\definecolor{termbg}{gray}{0.96}
\lstdefinestyle{term}{
  basicstyle=\ttfamily\small,
  backgroundcolor=\color{termbg},
  frame=single, rulecolor=\color{gray!40},
  breaklines=true, columns=fullflexible,
  xleftmargin=6pt, xrightmargin=6pt, framexleftmargin=6pt
}
\lstset{style=term}

\newcommand{\captura}[2]{%
  \begin{center}\fbox{\includegraphics[width=0.95\linewidth]{#1}}\end{center}%
  \begin{center}\small\textit{#2}\end{center}}
\newcommand{\ok}{\textcolor{green!50!black}{$\checkmark$}}

% ====== COMPLETAR ======
\newcommand{\estudiantes}{Randall Bolaños (carné \_\_\_\_\_\_\_\_) \and Javier Tenorio (carné \_\_\_\_\_\_\_\_)}
\newcommand{\fechaentrega}{\_\_\_ de octubre de 2026}
\newcommand{\enlacevideo}{\url{https://ENLACE-DEL-VIDEO}}
% =======================

\title{\textbf{SecureChat: comunicación segura con criptografía híbrida}\\\large Taller de Criptografía Aplicada --- Informe y documentación}
\author{\estudiantes}
\date{CE-1105 Seguridad de la Información \quad \fechaentrega}

\begin{document}
\maketitle
\tableofcontents
\bigskip

\noindent\textbf{Video de demostración:} \enlacevideo

\noindent\textbf{Código fuente:} \texttt{securechat.py}, \texttt{mitm\_proxy.py},
\texttt{test\_securechat.py} (ver \texttt{README.md} para instalar y ejecutar).

% ═══════════════════════════════════════════════════════════════
\section{Resumen de la solución}

SecureChat es una aplicación de consola en Python~3 que se ejecuta dos veces en la misma
computadora: una instancia \emph{escucha} (Bob) y otra se \emph{conecta} (Alice) mediante un
socket TCP real sobre \texttt{127.0.0.1}. Antes del primer mensaje ambas completan un
\emph{handshake} asimétrico (RSA-2048) que establece una clave de sesión AES-256 aleatoria sin
transmitirla nunca en claro; a partir de ahí cada mensaje se protege con AES-256-GCM, que cifra
y autentica. Se usa la biblioteca \texttt{cryptography} (versión 43).

Además se incluye \texttt{mitm\_proxy.py}, un proxy que se intercala entre ambos procesos,
imprime los bytes que realmente cruzan la red y puede alterar un byte del ciphertext para
demostrar la detección de manipulación, y un conjunto de 9 pruebas automáticas.

% ═══════════════════════════════════════════════════════════════
\section{Protocolo implementado}

\subsection{Algoritmos y tamaños de clave}

\begin{center}
\begin{tabularx}{\linewidth}{@{}lX@{}}
\toprule
\textbf{Elemento} & \textbf{Elección} \\ \midrule
Par asimétrico (cada participante) & RSA, $e=65537$, \textbf{2048 bits}. Privada en \texttt{claves/<nombre>\_privada.pem} con permisos 600. \\
Protección de la clave de sesión & \textbf{RSA-OAEP} con MGF1-SHA256 y SHA-256 (256 B de salida). \\
Autenticación del iniciador & Firma \textbf{RSA-PSS-SHA256} de Alice sobre (resumen del handshake $\|$ clave cifrada). \\
Clave de sesión & 32 bytes de \texttt{os.urandom} (CSPRNG), nueva en cada ejecución: \textbf{AES-256}. \\
Cifrado de mensajes & \textbf{AES-256-GCM}: nonce de 96 bits, tag de 128 bits. \\
Nonce & prefijo de dirección (4 B) $\|$ contador de 64 bits; nunca se repite. \\
Datos asociados (AAD) & \texttt{"SecureChat-v1"} $\|$ tipo de trama $\|$ $H$, con $H=\mathrm{SHA256}(\texttt{"SecureChat-v1"}\|\mathrm{HELLO}_A\|\mathrm{HELLO}_B)$. \\
\bottomrule
\end{tabularx}
\end{center}

\subsection{Handshake}

\begin{lstlisting}[basicstyle=\ttfamily\scriptsize,columns=fullflexible,keepspaces=true]
 Alice (conectar)                                       Bob (escuchar)
 HELLO {nombre, pubA}   ---------------------------------->   (publico)
                        <----------------------------------  HELLO {nombre, pubB}
 K = os.urandom(32)                   <- clave AES-256 de sesion
 C = RSA-OAEP(pubB, K)                <- solo privB la puede abrir
 S = RSA-PSS-Sign(privA, H || C)      <- prueba de que la envia Alice
 KEY_EXCHANGE  C || S   ---------------------------------->  Verify(pubA, S)
                                                             K = RSA-OAEP^-1(privB, C)
                        <----------------------------------  SESSION_OK = AES-GCM_K("SESSION-OK")
 verifica SESSION_OK
 ====== ambos poseen K  ->  "SESION SEGURA ESTABLECIDA"  ->  chat AES-256-GCM ======
\end{lstlisting}

La clave de sesión no está escrita en el código y solo cruza el socket cifrada con la clave
pública de Bob. Las claves privadas jamás se transmiten ni se imprimen. El mensaje
\texttt{SESSION\_OK} sirve de confirmación: al estar cifrado y autenticado con $K$, Alice
comprueba que Bob obtuvo exactamente la misma clave. De forma opcional,
\texttt{-{}-clave-par <pem>} fija la clave pública esperada del otro extremo (\emph{pinning}) para
impedir que un intermediario sustituya las claves públicas en el HELLO.

\subsection{Estructura de los mensajes transmitidos}

Toda trama tiene la forma \texttt{[longitud 4 B][tipo 1 B][payload]}:

\begin{center}
\begin{tabular}{@{}llll@{}}
\toprule
\textbf{Tipo} & \textbf{Nombre} & \textbf{Payload} & \textbf{¿Protegido?} \\ \midrule
\texttt{0x01} & HELLO        & JSON \{nombre, clave pública PEM\} & No (es público) \\
\texttt{0x02} & KEY\_EXCHANGE & RSA-OAEP($K$) 256 B $\|$ firma PSS 256 B & RSA \\
\texttt{0x03} & SESSION\_OK   & nonce $\|$ ciphertext $\|$ tag & AES-256-GCM \\
\texttt{0x10} & DATA         & nonce 12 B $\|$ ciphertext $n$ B $\|$ tag 16 B & AES-256-GCM \\
\texttt{0x11} & BYE          & nonce $\|$ ciphertext $\|$ tag & AES-256-GCM \\
\bottomrule
\end{tabular}
\end{center}

Ejemplo real: la trama del mensaje \texttt{"Hola Bob"} (41 bytes) tomada de la ejecución:

\begin{lstlisting}[basicstyle=\ttfamily\tiny,columns=fullflexible,keepspaces=true]
0000 0025 | 10 | 0000 0001 0000 0000 0000 0000 | 89e8 19b0 2b49 bb5a | f42e c623 7b1f 7be7 c272 013e 490a 77a8
longitud   tipo   nonce (direccion=1, ctr=0)     ciphertext (8 B)      tag GCM (16 B)
\end{lstlisting}

Los bytes ASCII de ``Hola Bob'' (\texttt{486f 6c61 2042 6f62}) no aparecen en la trama.

\textbf{Nonce.} Los 4 primeros bytes identifican la dirección (\texttt{00000001} Alice$\to$Bob,
\texttt{00000002} Bob$\to$Alice) y los 8 restantes son un contador que solo crece. Así ningún
extremo repite un nonce propio y los nonces de ambos extremos nunca colisionan aunque compartan
$K$. El receptor exige además contador estrictamente creciente (anti-\emph{replay}).

% ═══════════════════════════════════════════════════════════════
\section{Desarrollo y evidencias}

La ejecución de las capturas usó tres terminales: Bob (\texttt{escuchar}, puerto 5000), el
proxy (puerto 6000 $\to$ 5000, alterando la 7.ª trama DATA) y Alice (\texttt{conectar} al 6000).
Las salidas completas están en la carpeta \texttt{evidencia/}.

\subsection{Pasos 1--3: conexión, material asimétrico y handshake}

Cada instancia genera (o carga) su par RSA-2048, muestra solo la ruta y la huella de la clave
pública (nunca el contenido de la privada), establece la conexión TCP indicando con quién está
conectada y completa los 4 pasos del handshake antes de habilitar el chat.

\captura{capturas/01_handshake_bob.png}{Bob: claves, conexión TCP aceptada, handshake y ``sesión segura establecida''.}
\captura{capturas/02_handshake_alice.png}{Alice: la clave de sesión se genera en tiempo de ejecución y solo viaja cifrada con RSA-OAEP (256 B).}

\subsection{Pasos 4--5: mensajes protegidos y conversación bidireccional}

El emisor muestra el plaintext local, el nonce, el ciphertext, el tag y los bytes exactos
escritos en el socket; el receptor muestra los bytes recibidos y, \emph{solo tras verificar el
tag}, el plaintext recuperado. Se enviaron tres mensajes en cada dirección; el contador del
nonce avanza (0, 1, 2, \dots) en cada mensaje.

\captura{capturas/03_chat_alice.png}{Terminal de Alice: mensajes propios (plaintext $\to$ ciphertext) y mensajes de Bob (ciphertext $\to$ plaintext verificado).}
\captura{capturas/04_chat_bob.png}{Terminal de Bob: los bytes recibidos son idénticos a los enviados por Alice.}

\textbf{Lo que ve un tercero en la red.} El proxy imprime cada trama tal como cruza el socket.
Solo los HELLO son legibles (contienen claves \emph{públicas}); todo lo demás es ciphertext.

\captura{capturas/05_proxy_trafico.png}{Tráfico capturado por el proxy: ningún mensaje del chat aparece en claro.}

Sobre la captura binaria completa (\texttt{evidencia/captura.bin}):
\begin{lstlisting}
$ grep -a -c "Hola Bob"  captura.bin   ->  0
$ grep -a -c "4321"      captura.bin   ->  0
$ grep -a -c "PRIVATE"   captura.bin   ->  0
\end{lstlisting}

\subsection{Paso 6: detección de manipulación}

Se hicieron dos pruebas: (a) el proxy invirtió un bit del primer byte del ciphertext de la
trama ``Transferir 100 colones a Bob'' y (b) Bob usó \texttt{/alterar}, que modifica un byte del
ciphertext antes de enviarlo. En ambos casos el receptor detecta que el tag GCM no coincide,
\textbf{rechaza} el mensaje y no muestra ningún plaintext. El siguiente mensaje legítimo se
acepta con normalidad.

\captura{capturas/06_alteracion_rechazada_bob.png}{Bob rechaza la trama alterada por el proxy; después acepta ``Este si llega bien''.}
\captura{capturas/07_alteracion_rechazada_alice.png}{Alice rechaza el mensaje que Bob alteró con \texttt{/alterar}.}

\subsection{Pruebas automáticas}

\captura{capturas/08_pruebas_automaticas.png}{Las 9 pruebas de requisitos de seguridad pasan.}

% ═══════════════════════════════════════════════════════════════
\section{Verificación de requisitos mínimos}

\begin{center}
\begin{tabularx}{\linewidth}{@{}cX X@{}}
\toprule
 & \textbf{Requisito} & \textbf{Cómo se cumple} \\ \midrule
\ok & Python 3 y sockets TCP/IP & \texttt{socket.AF\_INET/SOCK\_STREAM} sobre 127.0.0.1 \\
\ok & Biblioteca reconocida & \texttt{cryptography} (hazmat RSA, AESGCM) \\
\ok & Handshake asimétrico & RSA-2048 OAEP + firma PSS \\
\ok & Clave de sesión dinámica & \texttt{os.urandom(32)} por sesión; prueba \emph{clave\_distinta\_en\_cada\_sesion} \\
\ok & Ninguna privada por el socket & solo se envían públicas; prueba de captura y \texttt{grep PRIVATE} = 0 \\
\ok & Cifrado autenticado & AES-256-GCM en DATA, SESSION\_OK y BYE \\
\ok & Nonce correcto y no repetido & prefijo de dirección + contador; prueba con 2000 cifrados \\
\ok & Rechazo de tag inválido & \texttt{InvalidTag} $\Rightarrow$ ``MENSAJE RECHAZADO''; prueba alterando cada byte \\
\ok & Evidencia visible & plaintext local, ciphertext enviado, ciphertext recibido, plaintext recuperado \\
\ok & No Base64/hex como cifrado & hex solo para mostrar bytes en pantalla \\
\ok & No ECB & GCM (modo contador autenticado); mismo texto $\Rightarrow$ ciphertexts distintos \\
\bottomrule
\end{tabularx}
\end{center}

% ═══════════════════════════════════════════════════════════════
\section{Preguntas de cierre}

\begin{enumerate}
\item \textbf{¿Qué problema resuelve la criptografía asimétrica durante el handshake?}
  El de la distribución de claves: Alice y Bob no comparten ningún secreto previo y el canal es
  observable. Con RSA-OAEP Alice cifra la clave de sesión con la clave \emph{pública} de Bob,
  que puede conocer cualquiera, y solo la \emph{privada} de Bob, que nunca sale de su equipo,
  puede recuperarla. Así acuerdan un secreto simétrico sin enviarlo nunca en claro. La firma
  RSA-PSS de Alice autentica además quién generó esa clave.

\item \textbf{¿Por qué la conversación utiliza criptografía simétrica después del handshake?}
  Porque AES-GCM es órdenes de magnitud más rápido que RSA, no tiene límite práctico de tamaño
  (RSA-OAEP-2048 solo cifra unos 190 B por operación y produce 256 B) y aporta autenticación
  integrada. RSA se usa una sola vez para acordar $K$, y la primitiva eficiente protege todos los
  mensajes: eso es la criptografía híbrida.

\item \textbf{¿Qué información puede observar un tercero que capture los paquetes?}
  IPs y puertos; el momento y el tamaño de cada trama, lo que deja inferir la longitud de cada
  mensaje; los HELLO, con los nombres y las claves \emph{públicas}; la clave de sesión cifrada
  con RSA y la firma; y, por cada mensaje, el tipo, el nonce, el ciphertext y el tag. No puede
  ver el contenido, la clave de sesión ni las claves privadas (ver captura del proxy).

\item \textbf{¿Por qué el nonce puede viajar junto al ciphertext y por qué no debe repetirse?}
  El nonce no es secreto: solo garantiza que cada cifrado sea único, y la seguridad descansa en
  la clave. El receptor lo necesita para descifrar. Pero GCM cifra en modo contador: repetir un
  nonce con la misma clave reutiliza el mismo flujo de cifrado ($C_1\oplus C_2 = P_1\oplus P_2$,
  lo que filtra el contenido) y permite recuperar la subclave de autenticación GHASH y
  \emph{falsificar} tags. Por eso usamos un contador monotónico con prefijo por dirección.

\item \textbf{¿Qué función cumple el authentication tag?}
  Es un código de autenticación de 128 bits que se calcula con $K$ sobre el ciphertext y los
  datos asociados. El receptor lo verifica \emph{antes} de entregar el plaintext: si cambió un
  solo bit del nonce, el ciphertext, el tag o el contexto, la verificación falla y el mensaje se
  rechaza. Garantiza integridad y autenticidad (lo envió alguien que posee $K$). Cifrar sin
  autenticar permitiría modificar mensajes sin que nadie lo detectara.

\item \textbf{¿Qué ocurriría si la clave de sesión se transmitiera sin protección?}
  Cualquiera que capturara el tráfico obtendría $K$ y podría descifrar toda la conversación y
  además fabricar o modificar mensajes con tags válidos, que el receptor aceptaría como
  auténticos. AES-GCM no aportaría ninguna seguridad. Codificarla en Base64 o hex no cambia
  nada: codificar no es cifrar.

\item \textbf{¿Qué ocurriría si un atacante obtiene una de las claves privadas?}
  Con la privada de \emph{Bob} podría descifrar el KEY\_EXCHANGE de cualquier sesión capturada,
  pasada o futura, obtener $K$ y leerla completa, porque el transporte de clave con RSA no
  ofrece \emph{forward secrecy}. También podría suplantar a Bob. Con la privada de
  \emph{Alice} podría firmar intercambios falsos y suplantarla ante Bob, aunque no descifrar
  sesiones pasadas. Las mitigaciones son proteger las privadas (permisos, contraseña, HSM),
  revocarlas y regenerarlas, y usar un intercambio efímero (ECDHE/X25519) firmado, como TLS~1.3,
  para lograr \emph{forward secrecy}.
\end{enumerate}

% ═══════════════════════════════════════════════════════════════
\section{Conclusiones y limitaciones}

La práctica separa con claridad los dos papeles. La criptografía asimétrica resuelve una sola
vez el problema de acordar un secreto sin canal seguro previo, y la simétrica autenticada
protege eficientemente cada mensaje. La prueba de manipulación muestra que cifrar no basta:
sin el tag, un atacante podría alterar el ciphertext sin ser detectado.

Limitaciones conscientes del diseño:
\begin{itemize}
  \item La autenticidad de las claves públicas depende de comparar huellas o de usar
        \texttt{-{}-clave-par}; no hay PKI.
  \item No hay \emph{forward secrecy}.
  \item La longitud de los mensajes es observable.
  \item Las privadas se guardan sin contraseña; para la demo solo se protegen con permisos 600.
\end{itemize}

\end{document}
